\section{Introduction}
\label{sec:intro}

Biomedical knowledge graphs (KGs) have become indispensable resources for drug discovery, disease understanding, and clinical decision support~\citep{Chandak2023, Huang2024}.
These graphs integrate diverse biomedical entities---genes, proteins, drugs, diseases, and biological processes---into rich relational structures that power downstream applications such as drug repurposing~\citep{Huang2024} and target identification.
However, biomedical knowledge is not static: new drug-disease associations, protein interactions, and gene-phenotype links are discovered continuously as databases like DisGeNET, the Human Phenotype Ontology (HPO), and Gene Ontology (GO) are updated.
A model trained on a KG snapshot from 2021 will encounter substantial new knowledge by 2023, raising the fundamental question: \emph{how can graph-based models continually learn from evolving biomedical KGs without catastrophically forgetting previously acquired knowledge?}

Continual graph learning (CGL) has emerged as an active research area addressing this challenge~\citep{Daruna2021}.
Methods such as Elastic Weight Consolidation (EWC)~\citep{Kirkpatrick2017}, experience replay~\citep{Rolnick2019}, and specialized frameworks like LKGE~\citep{Cui2023} have been proposed to mitigate catastrophic forgetting in knowledge graph embeddings.
However, existing CGL benchmarks suffer from a critical limitation: they construct temporal tasks through \emph{synthetic} partitioning of static, generic KGs.
For example, LKGE~\citep{Cui2023} partitions FB15k-237 and WN18RR into artificial time steps by randomly shuffling and splitting triples, while PS-CKGE~\citep{Wu2024} creates pseudo-temporal sequences from ICEWS~\citep{GarciaDuran2018} event data.
These approaches fail to capture the real dynamics of knowledge evolution---the addition of genuinely new entities and relations, the removal of deprecated knowledge, and the domain-specific patterns of change that characterize biomedical KGs.

In this paper, we introduce \ours (\textbf{Prime}KG for \textbf{C}ontinual \textbf{L}earning), the first CGL benchmark built on a real biomedical knowledge graph with genuine temporal evolution.
\ours is constructed from PrimeKG~\citep{Chandak2023}, a precision medicine KG integrating 20+ biomedical databases.
We reconstruct two temporal snapshots by querying the same upstream databases at different time points: $t_0$ (June 2021, the original PrimeKG release) and $t_1$ (July 2023, rebuilt from nine freely accessible databases).
The temporal difference between snapshots reveals substantial real-world evolution: 5.7M edges added, 889K edges removed, and 7.2M edges persisting---reflecting genuine changes in biomedical understanding rather than artificial partitioning.

Our benchmark evaluation of seven continual learning methods reveals several important findings.
First, our proposed Cross-Modal Knowledge Learner (CMKL), which leverages multimodal node features through cross-attention fusion and modality-aware EWC, achieves the highest average link prediction performance (AP\,=\,0.063), outperforming Joint Training by 3.5$\times$.
Second, EWC provides meaningful forgetting reduction (AF\,=\,0.017) compared to naive sequential training (AF\,=\,0.021), validating its utility in the biomedical domain.
Third, domain-specific challenges emerge: the gene/protein entity type dominates the graph (73\% of triples), creating a type imbalance that affects all methods, and the DistMult decoder~\citep{Yang2015} substantially outperforms TransE~\citep{Bordes2013} on this heterogeneous biomedical graph.

\vspace{0.3em}
\noindent Our contributions are as follows:
\begin{enumerate}
    \item \textbf{Benchmark.} We introduce \ours, the first CGL benchmark on a real biomedical KG with genuine temporal evolution from nine authoritative databases, comprising 129K+ nodes, 8.1M+ edges, 10 node types, 30 relation types, and 10 continual learning tasks.
    \item \textbf{Multimodal features.} We provide multimodal node representations---textual (BiomedBERT), molecular (Morgan fingerprints), and structural (R-GCN)---enabling the study of modality-specific forgetting patterns.
    \item \textbf{Comprehensive evaluation.} We benchmark seven continual learning methods across three evaluation tracks (link prediction, KGQA, node classification) with five random seeds and report six CL metrics (AP, AF, BWT, FWT, REM, MRR).
    \item \textbf{Open release.} We release the benchmark data, construction pipeline, all baseline implementations, and trained models to facilitate reproducible research.
\end{enumerate}

\begin{table}[t]
\centering
\caption{Comparison of \ours with existing continual graph learning benchmarks. \ours is the first to offer real temporal evolution on a biomedical KG with multimodal features and multiple evaluation tasks.}
\label{tab:comparison}
\small
\begin{tabular}{lccccccc}
\toprule
\textbf{Benchmark} & \textbf{Domain} & \textbf{Real Temp.} & \textbf{Multimodal} & \textbf{\#Entities} & \textbf{\#Relations} & \textbf{\#Tasks} & \textbf{Eval Tasks} \\
\midrule
LKGE~\citep{Cui2023} & General & \xmark & \xmark & 15K & 237 & 5 & LP \\
PS-CKGE~\citep{Wu2024} & General & \xmark & \xmark & 15K & 237 & 5 & LP \\
ICEWS~\citep{GarciaDuran2018} & Political & \cmark & \xmark & 13K & 256 & -- & LP \\
\midrule
\textbf{\ours (ours)} & \textbf{Biomedical} & \cmark & \cmark & \textbf{129K+} & \textbf{30} & \textbf{10} & \textbf{LP, KGQA, NC} \\
\bottomrule
\end{tabular}
\end{table}

